
\section{Introduction}
\label{sec:intro}


Multidimensional point cloud data sets are becoming pervasive in numerous agricultural applications. Advances in phenotyping technologies that measure various crop attributes, coupled with an increased adoption of field sensing for environmental monitoring (e.g., temperature, humidity, soil moisture) have led to increased availability of multidimensional data sets in agriculture.
While most of these variables are temporal, some variables may also encode static attributes that describe spatial locators or the crop varieties (cultivars) grown at those locations.

Given such a complex spatiotemporal data set, we are interested in understanding how different cultivars (with different genotypes G) respond to different environmental factors (E) to affect their phenotypes (P) \cite{tardieu_plant_2017}. This question of decoding the G$\times$E$\rightarrow$ P interaction is at the center of modern-day phenomics.
However, prior to understanding this complex interaction, historical data sets offer a more immediate opportunity to understand how different genotypes relate to one another by their phenotypic behavior. For instance, extracting different patterns of phenotypic behavior---be it conserved or divergent---among various subgroups of cultivars could help classify cultivars into  behavioral groups. Class information could help in subsequent prediction tasks associated with those cultivars. However, in agricultural data sets, such a classification may not be trivially observable from data, particularly for temporally changing phenotypes.  Given the complex nature of traits, a subset of cultivars that show similar phenotypic behavior during one part of the season may possibly diverge at other times.  Secondly, even within a single cultivar, there could be  variability observed across the different seasons, as phenotypic plasticity is a well-established phenomenon in plants \cite{schlichting1986evolution}.
Consequently, it is important for any downstream machine learning workflow to incorporate these complex structural relationships among cultivars in order to improve the efficacy of the prediction tasks. 

In this paper, we model the problems of extracting cross-cultivar  relationships and intra-cultivar variability patterns as one of a structure discovery process. More specifically, given a multi-dimensional point cloud, where each point represents a phenotypic observation of a cultivar in time, we model the problem of identifying structural patterns in data using topological data analysis. In particular, we explore the use of persistent homology \cite{edelsbrunner_persistent_2013}, an active research area within the field of applied algebraic  topology. This is a branch of mathematics that studies \emph{shapes} of data and spaces using algebraic techniques. Topology works with coordinate-free representation of shapes (simplicial complexes) \cite{munkres2018elements},
which are also more robust to small changes in data or to missing data \cite{lum_extracting_2013}.

There are two major classes of techniques within applied  topology---\texttt{mapper} \cite{singh_topological_2007} and persistent homology \cite{edelsbrunner_persistent_2013}. Here, we explore persistent homology to study point clouds as it is better equipped to elucidate topological features that persist over multiple scales (including temporal scales).
Although persistent homology has been widely used in a number of other application domains, it is yet to be explored in any serious depth for agricultural data sets. 
To the best of our knowledge, it has only been applied to decode plant leaf shapes \cite{zhang2021mfcis}. However, the focus of this paper is different; we aim to identify patterns between cultivars and within cultivars based on phenotypic behavior. 

As a concrete application case study, we study the cold hardiness of multiple grape cultivars.
Cold hardiness is a trait that measures how resilient a variety is to cold temperatures. Specialty crops such as grapes and apples can incur a significant loss when the air temperature drops below certain cold hardiness thresholds \cite{mills_cold-hardiness_2006}. However, these thresholds are not fixed, change by the time of the season, and vary by the cultivars. 
Furthermore, due to a large number of cultivars and their divergent behavior across different growing conditions, it becomes important to study: a) the relationships between different cultivars by their cold hardiness trait, and b) any trait level variability as seen in the same cultivar across different seasons. 
Elucidating these relationships will help field scientists devise better frost/cold mitigation protocols customized and effective when applied to different cultivars. 
They could also help us improve the precision of current state-of-the-art cold hardiness prediction models (e.g., \cite{ferguson_modeling_2014}).

%\delete{OLD TEXT from here on}


%In this work, we apply topological data analysis (TDA), a class of techniques from computational algebraic topology. TDA can help in discovering patterns (conserved or discrepant) from high-dimensional, possibly incomplete, or noisy point clouds \cite{lum_extracting_2013}. Using TDA, we implemented a computational framework to analyze such data in a manner that is insensitive to the particular performance metric chosen. The two approaches commonly used in topological data analysis are the mapper algorithm \cite{singh_topological_2007} and persistence homology \cite{edelsbrunner_persistent_2013}. The mapper algorithm allows you to transform a point cloud into a graph, enabling you to examine branching events that signify differences between cultivars. On the other hand, persistence homology is used to discern topological features of data such as connected components, holes, graph structure, etc. Using these two methodologies, we found interesting correlations between cultivars and across different seasons for the same cultivar.
